\documentclass[12 pt, letterpaper]{hmcpset}
\usepackage{hw-preamble}

\name{Jayden Thadani}
\class{MATH 131 --- Mathematical Analysis I}
\assignment{Homework 8}
\duedate{9th November 2024}

\begin{document}

\begin{problem}[Rudin chapter 4, exercise 4]
	Let $f$ and $g$ be continuous mappings of a metric space $X$ into a metric space $Y$ and let $E$ be a dense subset of $X$. Prove that $f^{\text{img}}(E)$ is dense in $f^{\text{img}}(X)$. If $g(p) = f(p)$ for all $p \in E$, prove that $g(p) = f(p)$ for all $p \in X$.
\end{problem}

\begin{solution}
	Since $f$ is continuous, given any $\varepsilon > 0$ and $f(x) \in f^{\text{img}}(X)$ there exists a $\delta > 0$ so that the open ball around $x$ is mapped into the open ball around $f(x)$ --- $f^{\text{img}}(B_{\delta}(x)) \subset B_{\varepsilon}(f(x))$. Since $E$ is dense in $X$, there is some $p \in E \cap B_{\delta}(x)$, and thus, $f(p) \in f^{\text{img}}(E) \cap B_{\varepsilon}(f(x))$. That is, for any $\varepsilon > 0$, and any $f(x) \in f^{\text{img}}(X)$, the intersection $f^{\text{img}}(E) \cap B_{\varepsilon}(f(x))$ is non-empty --- $f^{\text{img}}(E)$ is dense in $f^{\text{img}}(X)$.

Suppose $f$ and $g$ agree on the dense subset $E$. Since $E$ is dense in $X$, for any point $x \in X$, there is a sequence $x_{n} \in E$ with $x_{n} \to x$. Since $f$ and $g$ are continuous, the $f(x_{n}) \to f(x)$ and the $g(x_{n}) \to g(x)$. Since $f(x_{n}) = g(x_{n})$, we have $f(x) = g(x)$.
\end{solution}

\pagebreak

\begin{problem}[Rudin chapter 4, exercise 6]
	If $f$ is defined on $E$, the \emph{graph} of $f$ is the set of points $(x, f(x))$ for $x \in E$. In particular, if $E$ is the set of real numbers, and $f$ is real-valued, the graph of $f$ is a subset of the plane.

	Suppose $E$ is compact, and prove that $f$ is continuous on $E$ if and only if its graph is compact.
\end{problem}

\begin{solution}
	First, suppose $f$ is continuous. Then we have that $f^{\text{img}}(E)$ is compact. Since $E$ and $f^{\text{img}}(E)$ are bounded, it follows that $E \times f^{\text{img}}(E)$ is bounded, and thus, that the graph of $f$, which is a subset of $E \times f^{\text{img}}(E)$ is bounded too. Thus, we only need to show that the graph of $f$ is closed. 

	Suppose $(x_{0}, y_{0})$ is a limit point of the graph of $f$. Thus, for any $\varepsilon > 0$, we can find $(x, f(x))$ within distance $\varepsilon$ of $(x_{0}, y_{0})$. Thus, for  $\delta = \varepsilon$, we can force  $\lvert f(x) - y_{0} \rvert < \varepsilon$ for all $\lvert x - x_{0} \rvert < \delta$. That is, $\lim_{x \to x_{0}} f(x) = y_{0}$. Thus, by continuity, $f(x_{0}) = y_{0}$, and $(x_{0}, y_{0})$ is in the graph of $f$. Thus, the graph of $f$ is closed, and since it is bounded, it is also compact.

	Now we prove the converse. Suppose that the graph of $f$ is compact and $\lim_{x \to x_{0}} f(x) = y_{0}$. Then $(x_{0}, y_{0})$ is a limit point of the graph of $f$ since for every, $\varepsilon > 0$ there is $\delta > 0$ with $0 < \lvert x - x_{0} \rvert < \delta$ forcing $\lvert f(x) - y_{0} \rvert < \varepsilon$. Thus, for every $r > 0$, we can choose $\varepsilon < r$ so that we can choose some $x$ with $\lvert x - x_{0} \rvert < \min \{ \delta, \sqrt{r^{2} - \varepsilon^{2}} \}$, and thus, $(x, f(x))$ is within $r$ of $(x_{0}, y_{0})$. That is, $(x_{0}, y_{0})$ is a limit point of the graph of $f$. Thus, $(x_{0}, y_{0})$ is a point of the graph of $f$, giving us $f(x_{0}) = y_{0}$. That is, at every point, $f$ equals its limit --- $f$ is continuous.
\end{solution}

\pagebreak

\begin{problem}[Rudin chapter 4, exercise 7]
	If $E \subset X$ and if $f$ is a function defined on $X$, the \emph{restriction} of $f$ to $E$ is the function $g$ whose domain of definition is  $E$ such that $g(p) = f(p)$ for $p \in E$.

	Define $f$ and $g$ on $\mathbb{R}^{2}$ by $f(0, 0) = g(0, 0) = 0$, $f(x, y) = xy^{2} / (x^{2} + y^{4})$, $g(x, y) = x y^{2}/(x^{2} + y^{6})$ if $(x, y) \neq (0, 0)$. Prove that $f$ is bounded on $\mathbb{R}^{2}$, that $g$ is unbounded in every neighbourhood of $(0, 0)$ and that $f$ is not continuous at $(0, 0)$; nevertheless the restrictions of both $f$ and $g$ to every straight line in $\mathbb{R}^{2}$ are continuous.
\end{problem}

\begin{solution}
	First we show that $f$ is bounded. We know that for any $x, y \in \mathbb{R}$, $(x - y^{2})^{2} \ge 0$. Thus,
	\[
		x^{2} + y^{4} - 2 x y^{2} \ge 0.
	\]
	Then
	\[
		x^{2} + y^{4} \ge 2 x y^{2}
	\]
	and finally,
	\[
		\frac{x y^{2}}{x^{2} + y^{4}} \le \frac{1}{2}.
	\]
	Thus, $f$ is bounded.

	All straight lines apart from $x = 0$ are given by $y = \lambda x$ for some $\lambda \in \mathbb{R}$. Thus, the restriction of $f$ to the line is $f_{\lambda}$ given by
	\[
		f_{\lambda}(x) = \frac{\lambda^{2} x^{3}}{x^{2} + \lambda^{4} x^{4}}
	\]
	for $x \neq 0$ and $f(0) = 0$. Since $f_{\lambda}$ is the quotient of non-zero functions everywhere except $x = 0$, we just have to verify $\lim_{x \to 0} f_{\lambda}(x) = 0$.
	\begin{equation*}
		\lim_{x \to 0} \frac{\lambda^{2} x^{3}}{x^{2} + \lambda^{4} x^{4}} = \lim_{x \to 0} \frac{\lambda x}{1 + \lambda^{4} x^{2}} = 0.
	\end{equation*}
	For the line $x = 0$, the restriction is $0$ everywhere, and thus, the restriction is continuous

	Similarly we can define the restriction $g_{\lambda}$ defined by
	\[
		g_{\lambda}(x) = \frac{\lambda x^{3}}{x^{2} + \lambda^{6} x^{6}}
	\]
	for $x \neq 0$ and $g(0) = 0$. For the same reasons as previously, we only need to show that $\lim_{x \to 0} g_{\lambda}(x) = 0$. This follows as
	\[
		\lim_{x \to 0} \frac{\lambda^{2} x^{3}}{x^{2} + \lambda^{6} x^{6}} = \lim_{x \to 0} \frac{\lambda^{2} x}{1 + \lambda^{6} x^{4}} = 0.
	\]
	and on the line $x = 0$, the restriction is $0$ everywhere, and thus, the restriction is continuous.
\end{solution}

\pagebreak

\begin{problem}[Rudin chapter 4, exercise 8]
	Let $f$ be a real uniformly continuous function on the bounded set $E$ in $\mathbb{R}$. Prove that $f$ is bounded on $E$.

	Show that the conclusion is false if boundedness of $E$ is omitted from the hypothesis.
\end{problem}

\begin{solution}
	Suppose $f$ is uniformly continuous on $E$. Then $f$ is the union of finitely many connected components --- intervals. Choose some $\varepsilon > 0$, and find $\delta > 0$ such that for all $x_{0} \in E$, we have $\lvert f(x) - f(x_{0}) \rvert < \varepsilon$ for $\lvert x - x_{0} \rvert < \delta$. Then on each connected component $(a, b)$ of $E$, choose $[a + \delta, b - \delta]$. Since this is compact, so is its image. Applying uniform continuity to the $a$ and $b$, we get that $f$ is bounded on all of $(a, b)$. Thus, there are finitely many bounds on $f$ and $f$ is bounded. (Similarly if the connected component is $(a, b]$, we consider $[a + \delta, b]$, for $[a, b)$ we consider $[a, b - \delta]$, and $[a, b]$ is compact itself).

	Consider $f : x \mapsto x$ which is uniformly continuous on $\mathbb{R}$. $f$ is unbounded.
\end{solution}	

\pagebreak

\begin{problem}[Rudin chapter 4, exercise 11]
	Suppose $f$ is a uniformly continuous mapping of a metric space $X$ into a metric space $Y$ and prove that $\{ f(x_{n}) \}_{n}$ is a Cauchy sequence in $Y$ for every Cauchy sequence $\{ x_{n} \}_{n}$ in $X$. Use this result to give an alternative proof of the theorem stated in exercise 13.
\end{problem}

\begin{solution}
	Consider the Cauchy sequence $x_{n}$. Since $f$ is uniformly continuous, we have that for any $\varepsilon > 0$, there is a uniform $\delta > 0$ such that at any $x_{n}$ we can force $d(f(x), f(x_{n})) < \varepsilon$ for any $d(x, x_{n}) < \delta$. But then, for any $x_{m}, x_{n}$ with $d(x_{n} - x_{m}) < \delta$, we have $d(f(x_{n}), f(x_{m})) < \varepsilon$.

	Since $x_{n}$ is a Cauchy sequence, there exists a positive integer $N$ such that for all $m, n > N$ we have $d(x_{n}, x_{m}) < \delta$ for any positive $\delta$. Now it follows that for any $\varepsilon > 0$ there exists a positive integer $N$ such that for all $m, n > N$, we have $d(f(x_{n}), f(x_{m}))  < \varepsilon$. That is, $f(x_{n})$ is a Cauchy sequence.

	The theorem in exercise 13 asks us to show that a uniformly continuous real function $f : E \to \mathbb{R}$ from a dense subset $E \subset X$ has a continuous extension to all of $X$. Since $E$ is dense in $X$, for every $x \in X$, there exists a sequence $x_{n} \in E$ converging to $X$. Since this sequence converges, it is Cauchy. By our previous result, the sequence $f(x_{n})$ must also be Cauchy as well. Since $\mathbb{R}$ is complete, this sequence converges. Thus, define the extension
	\[
		g(x) = \lim_{n \to \infty} f(x_{n}) \quad \text{for} \quad x_{n} \to x.
	\]
	Note that this is clearly equal to the continuous $f$ on $E$, but we need to show that this is well defined. 

	For any two sequences $x_{n}, z_{n} \to x$, note that for sufficiently large $n$, we can get $d(x_{n}, z_{n}) < \delta$ for any $\delta > 0$. (This follows by applying the triangle inequality to the convergence of both sequences to $x$).  Thus, by uniform continuity, for sufficiently large $n$, we can get $\lvert f(x_{n}) - f(z_{n}) \rvert < \varepsilon$. Since both of these are convergent sequences, it follows by the triangle inequality that since both of them converge to the same real number.

	This serves as a proof that $g$ is well defined, and since $E$ is dense in $X$, any sequence converging to $x$ is arbitrarily close to a sequence in $E$, and so this also serves as a proof of sequential continuity. That is, $g$ is a continuous extension of $f$.
\end{solution}

\pagebreak

\begin{problem}[Problem 6]
	Prove that there is no continuous bijection between the interval $[0, 1]$ and the triangle in $\{ (x, 0) \mid x \in [0, 1] \} \cup \{ (x, x) \mid x \in [0, 1] \} \cup \{ (1, y) \mid y \in [0, 1] \}$.
\end{problem}

\begin{solution}
	Denote the triangle by $\Delta$. Note since both, the triangle and the interval are compact, a continuous bijection between them has a continuous inverse. Thus, without loss of generality, we can suppose for contradiction that there was a continuous bijection $f : \Delta \to [0, 1]$. 

	Pick some $a \in (0, 1)$. Since $f$ is a bijection, there is a unique $f^{-1}(a)$. Consider the restriction of $f$ to $\Delta \setminus \{ f^{-1}(a) \}$. This function has range $[0, 1] \setminus \{ a \}$ which is not connected. However, the restriction is still continuous since it is still equal to its limit everywhere in $\Delta \setminus \{ f^{-1}(a) \}$, and $\Delta \setminus \{ f^{-1}(a) \}$ is connected (since one can see that it is clearly path connected). This is a contradiction since the image of a connected set under a continuous function is connected. Thus, there is no continuous bijection $f : \Delta \to [0, 1]$.
\end{solution}

\end{document}
